\setcounter{chapter}{1}
\setcounter{section}{0}
\counterwithout{section}{chapter}
\chapter{Les systèmes de numération : Bases et changement de bases}


On appelle base d'un système de numération de position, le nombre de chiffres distincts permettant l'écriture de n'importe quel nombre entier naturel. Grâce au chapitre précédent, vous savez maintenant compter en base 10, 16 et 2. Nous allons généraliser ce principe pour toutes bases et nous verrons comment passer d'une base à une autre.\\

\section{Base d'un système de numération}

\subsection{Décomposition décimale}
La numération décimale utilise 10 chiffres : 0, 1, 2, 3, 4, 5, 6, 7, 8 et 9.\\
L'écriture du nombre peut se traduire avec l'égalité suivante $275_{(10)}$:
\[275= 2 \times 100 + 7 \times 10 + 5 \times 1= 2\times10^2 + 7\times10^1+ 5\times10^0\]
On dit que $275$ est la représentation en base 10 appelée écriture décimale.\\

On peut s'aider d'un tableau :
\begin{tabular}{lccccc}
& $10^4$ & $10^3$ & $10^2$ & $10^1$ & $10^0$\\
& $10000$ & $1000$ & $100$ & $10$ & $1$\\
\hline
$275_{(10)} = $ & 0 & 0 & 2 & 7 & 5\\
$56_{(10)} = $ & 0 & 0 & 0 & 5 & 6\\
$12_{(10)} = $ & 0 & 0 & 0 & 1 & 2\\

\end{tabular}\\
On en déduit : \\
$275= 2 \times 100 + 7 \times 10 + 5 \times 1= 2\times10^2 + 7\times10^1+ 5\times10^0$\\
$56_{(10)} = 5\times 10^1 + 6\times 10^0$\\
$12_{(10)} = 1\times 10^1 + 2\times 10^0$\\
\begin{Avous}
    Écrire une égalité semblable pour les nombres 1 234 et 208.\\[4cm]
\end{Avous}

\subsection{Numération hexadécimale et conversion en décimal}
La numération hexadécimale utilise 16 chiffres : 0, 1, 2, 3, 4, 5, 6, 7, 8, 9, A, B, C, D, E et F.\\
A représente donc le nombre $10$, B le nombre $11$ etc ...\\
\begin{center}
    

\begin{tabular}{lccc}
    & $16^2$ & $16^1$ & $16^0$\\
    & $256$ & $16$ & $1$\\
    \hline
    $38_{(16)} = $ & 0 & 3 & 8\\
    $B_{(16)} = $ & 0 & 0 & B\\
    \end{tabular}\\
\end{center}

    On déduit : \\
    $38_{(16)} = 3\times 16^1 + 8\times 16^0$\\
    $B_{(16)} = 11\times 16^0$\\

    
En réalité 11 est déjà la conversion du chiffre B hexadécimal en décimal. Si vous résolvez le calcul, on obtient :\\
    $38_{(16)} = 3\times 16^1 + 8\times 16^0 = 56_{(10)}$\\
    $B_{(16)} = 11\times 16^0 = 11_{(10)}$\\

    
Grâce à cette méthode on peut rapidement dire que la conversion du nombre hexadécimal 38 vaut 56 en décimal.
\begin{Avous}
     Traduire en nombre décimal $B8C_{(16)}$ et $FF_{(16)}$\\[3cm]
\end{Avous}


\subsection{Numération binaire et conversion en décimal}

C'est donc le système de base 2 qu'utilisent les ordinateurs. Les chiffres sont au nombre de 2 : 0 et 1. Ainsi le nombre $12 = 1100_{(2)}$ est stocké par $4$ circuits mémoires respectivement dans les états $1, 1, 0, 0$.\\
\begin{center}
\begin{tabular}{lcccccccc}
    & $2^7$ & $2^6$ & $2^5$ & $2^4$ & $2^3$ & $2^2$ & $2^1$ & $2^0$\\
    & $128$ & $64$ & $32$ & $16$ & $8$ & $4$ & $2$ & $1$\\
    \hline
    $1100_{(2)} = $ & 0 & 0 & 0 & 0 & 1 & 1 & 0 & 0\\
    $111000_{(2)} = $ & 0 & 0 & 1 & 1 & 1 & 0 & 0 & 0\\
    \end{tabular}
\end{center}

    On déduit : \\
    $ 111000_{(2)} = 1\times 2^5 + 1\times 2^4 + 1\times 2^3 + 0\times 2^2 + 0\times 2^1 + 0\times 2^0 = 56_{(10)}$\\
    $1100_{(2)}= 1\times 2^3 + 1\times 2^2 + 0\times 2^1 + 0\times 2^0 = 12_{(10)} $\\

\begin{Avous}
     Traduire en nombre décimal $101010_{(2)}$ et $11011_{(2)}$\\[2.5cm]
\end{Avous}

\begin{Comprendre}
    Pour rappel, un nombre binaire de huit chiffres est un \textbf{octet}.\\
Quel est le plus grand nombre décimal que l'on peut représenter avec un octet ?\\[2cm]
\end{Comprendre}
%%%%%%%%%%%%%%%%%%%%%%%%%%%%%%%%%%%%%%%%%%%%%%%%%%%%%%%%%%ù
%%%%%%%%%%%%%%%%%%%%%%%%%%%%%%%%%%%%%%%%%%%%%%%%%%%%%%%%%%%%
% LA DEBUT
%%%%%%%%%%%%%%%%%%%%%%%%%%%%%%%%%%%%%%%%%%%%%%%%%%%%%%%%%%ù
%%%%%%%%%%%%%%%%%%%%%%%%%%%%%%%%%%%%%%%%%%%%%%%%%%%%%%%%%%%%


% \subection{Généralisation du changement de base}

% Dans un système de base b, un nombre sera noté :
% \begin{equation*}
%     (a_n a_{n-1} ... a_1 a_0 )_b = a_0 \times b^0 + a_1 \times b^1 + ... + a_{n-1} \times b^*{n-1} + a_n \times b^n
% \end{equation*}

% où $a_i$ est un chiffre compris entre $0$ et $b-1$.

%\subsection*{Généralisation du passage de la base b à la base 10}

% Passer de l'écriture en base b d'un nombre à son écriture décimale est aisé, il suffit de faire la somme de la multiplication de chaque chiffre d'indice k par son poids correspondant $b^k$ (b étant la base d'origine).\\

% On prend un exemple d'un nombre en base 16 vers un nombre binaire :
% Au préalable, on prendra soin de transcoder les lettres A à F en décimal (de 10 à 15).
% Ainsi : $FA9_{(16)}= 15 \times 16^2 + 10 \times 16^1 + 9 \times 16^0 = 3840 + 160 + 9 = 4009_{(10)}$

% \begin{Avous}
%     Essayez maintenant d'une base 8 vers la base 10. Donner l'écriture décimale de $76401_{(8)}$\\[2cm]
% \end{Avous}

\section{Conversion de la base 10 à une base $b$}

Pour écrire les entiers naturels en base $b$, on a besoin de $b$ chiffres. Quand on a $n$ objets, on les groupe par paquets de $b$, qu'on regroupe à leur tour en paquets de $b$ paquets, etc. Autrement dit, on fait une succession de divisions par $b$, jusqu'à obtenir un quotient égal à 1.

\subsection*{Conversion de la base 10 à la base 2}
Dans le cas d'une conversion en binaire, on va faire une succession de divisions par 2.
Un exemple pour le nombre $77_{(10)}$\\
\begin{center}
\begin{tabular}{rcl}
$77$ & $=$ & $2\times 38 + \mathbf{1}$\\
$38$ & $=$ & $2\times 19 + \mathbf{0}$\\
$19$ & $=$ & $2\times 9 + \mathbf{1}$\\
$9$ & $=$ & $2\times 4 + \mathbf{1}$\\
$4$ & $=$ & $2\times 2 + \mathbf{0}$\\
$2$ & $=$ & $2\times 1 + \mathbf{0}$\\
$1$ & $=$ & $2\times 0 + \mathbf{1}$\\
\end{tabular}
\end{center}
Le résultat sera donc $77_{(10)} = 1001101_{(2)}$ (lecture de bas vers le haut des restes)\\

Aux choix, on peut également utiliser la méthode du tableau. Reprenons les tableaux avec l'exemple du nombre décimal $77_{(10)}$ qu'on veut convertir en base 2 :
\begin{center}
\begin{tabular}{lcccccccc}
    & $2^7$ & $2^6$ & $2^5$ & $2^4$ & $2^3$ & $2^2$ & $2^1$ & $2^0$\\
    & $128$ & $64$ & $32$ & $16$ & $8$ & $4$ & $2$ & $1$\\
    \hline
    $77_{(10)} = $ &  &  &  &  &  &  &  & \\
    \end{tabular}
\end{center}

\begin{itemize}
    \item Etape 1 : on écrit le tableau avec toutes les puissances de 2. Cela va nous permettre de suivre le nombre de "paquets de 2".
    \item Etape 2 : Dans le nombre 77, puis-je 1 fois faire un paquet de 128 ? Ici non, donc j'écris 0 dans la case correspondante.
    \item Etape 3 : Dans le nombre 77, puis-je 1 fois faire un paquet de 64 ? oui, donc j'écris 1 dans la case correspondante. Puis j'enlève ce paquet de 64 : $77-64 = 13$
    \item Etape 4 : Dans le nombre 13, puis-je une fois faire un paquet de 32 ? non, donc j'écris 0
    \item Etape 5 : Dans le nombre 13, puis-je une fois faire un paquet de 16 ? non, donc j'écris 0
    \item Etape 6 : Dans le nombre 13, puis-je une fois faire un paquet de 8 ? oui, donc j'écris 1. Puis j'enlève le paquet de 8 : $13-8 = 5$
    \item Etape 7 : Dans le nombre 5, puis-je une fois faire un paquet de 4 ? oui, donc j'écris 1. Puis j'enlève le paquet de 4 : $5-4 = 1$
    \item Etape 8 : Dans le nombre 1, puis-je une fois faire un paquet de 2 ? non, donc j'écris 0.
    \item Etape 9 : Dans le nombre 1, puis-je une fois faire un paquet de 1? oui, donc j'écris 1 :  Puis j'enlève le paquet de 1 : $1-1 = 0$. Je suis à 0 donc j'ai terminé, s'il reste des colonnes à remplir, j'écris 0 dedans.
\end{itemize}

\begin{Avous}
    Convertir $140_{(10)}$ en binaire. 
    \begin{center}
\begin{tabular}{lcccccccc}
    & $2^7$ & $2^6$ & $2^5$ & $2^4$ & $2^3$ & $2^2$ & $2^1$ & $2^0$\\
    & $128$ & $64$ & $32$ & $16$ & $8$ & $4$ & $2$ & $1$\\
    \hline
    $140_{(10)} = $ &  &  &  &  &  &  &  & \\
    \end{tabular}
\end{center}
\vspace{3cm}
\end{Avous}

\subsection{Conversion de la base 10 à la base 16}
Si on veut de l'hexadécimal, on va diviser par 16. Exemple pour $335_{(10)}$ :\\
\begin{center}
\begin{tabular}{rcl}
$335$ & $=$ & $16\times 20 + \mathbf{15}$ \quad (15 s'écrit F)\\
$20$ & $=$ & $16\times 1 + \mathbf{4}$\\
$1$ & $=$ & $16\times 0 + \mathbf{1}$\\
\end{tabular}
\end{center}
Donc $335_{(10)} = 14F_{(16)}$.\\

Comme précédemment, aux choix, on peut également utiliser la méthode du tableau. Reprenons les tableaux avec l'exemple du nombre décimal $335_{(10)}$ qu'on veut convertir en base 16 :
\begin{Avous}

\begin{center}
\begin{tabular}{lccc}
    & $16^2$ & $16^1$ & $16^0$\\
    & $256$ & $16$ & $1$\\
    \hline
    $335_{(10)} = $ &  &  & \\
    \end{tabular}
\end{center}
    \vspace{3cm}
    
\end{Avous}

\subsection{Conversion de la base 2 à la base 16}
Pour passer de la base 2 à la base 16 :
\begin{itemize}
    \item On décompose ce nombre par tranches de 4 bits à partir du bit de poids faible/tout à droite
    \item On complète la dernière tranche par des 0 s'il y a lieu.
    \item On convertit chaque tranche en son symbole de la base 16.
    \item On réécrit à sa place le nouveau symbole par changements successifs de chaque groupe de 4 bits.
\end{itemize}

\newpage
Par exemple $11~1101_{(2)}$
\begin{center}
\begin{tabular}{lcccccr}
    & $2^3$ & $2^2$ & $2^1$ & $2^0$\\
    & $8$ & $4$ & $2$ & $1$ &\\
    \hline
    $1101_{(2)} = $ & 1 & 1 & 0 & 1 & $=8+4+1=13_{10}=D_{16}$\\
    $11_{(2)} = $  & 0 & 0 & 1 & 1& $=2+1=3_{10}=3_{16}$\\
    \end{tabular}
\end{center}
Ainsi $11~1101_{(2)} = 3D_{(16)}$
\\[-1cm]
\begin{Comprendre}
    Le nombre maximal qu'on peut écrire sur 4 bits est $1111_{(2)}$, soit $15_{(10)}$, soit $F_{(16)}$. Finalement, l'hexadécimal est presque un manière plus rapide d'écrire le binaire. C'est pour cela qu'il est fréquemment utilisé en informatique. Il est permet un compromis entre le code binaire des machines et une base de numération pratique à utiliser pour les ingénieurs (chaque chiffre correspond exactement à 4 chiffres binaires - 4 bits).
\end{Comprendre}

\subsection{Conversion de la base 16 à la base 2}
Pour passer de la base 16 à la base 2, cette conversion est l'opération inverse de la précédente.
\begin{itemize}
    \item On convertit chaque symbole hexadécimal du nombre en son écriture binaire (nécessitant au max 4 bits)
    \item Pour chaque tranche de 4 bits, on complète les bits de poids fort par des 0 s'il y a lieu.
    \item On regroupe toutes les tranches de 4 bits à partir du bit de poids faible, sous forme d'un seul nombre binaire.\\
\end{itemize}


\begin{Avous}
Convertir $1100011010_{(2)}$ en base 16 et $7CF21_{(16)}$ en base 2.\\
    \vspace{8cm}
\end{Avous}


\centering
{\huge{
Bravo, vous êtes maintenant incollables sur les conversions entre les systèmes de numération en base 2, 10 et 16 !}}

% astuce à ajouter il est fréquemment utilisé en informatique car il est permet un compromis entre le code binaire des machines et une base de numération pratique à utiliser pour les ingénieurs (chaque chiffre correspond exactement à 4 chiffres binaires - 4 bits).

% \begin{Avous}
%     Trouver la représentation en base 16 puis en base 2 du nombre 6725 (écrits en base 10):\\[2cm]
% \end{Avous}




% \section{Contrôle}
% Comptez jusqu'à $1010_{(2)}$ en binaire.
% Comptez jusqu'à $F0_{(16)}$ en hexadécimal.
% De combien de bits est composé un octet ?
% Quel est le plus grand nombre qu'on peut représenter avec un octet ?
% Combien fait $FF_{(16)}+1_{(16)}$ ?
% \begin{tabular}{lcccccccc}
%     & $2^7$ & $2^6$ & $2^5$ & $2^4$ & $2^3$ & $2^2$ & $2^1$ & $2^0$\\
%     & $128$ & $64$ & $32$ & $16$ & $8$ & $4$ & $2$ & $1$\\
%     \hline
%      & 0 & 0 & 0 & 1 & 1 & 1 & 0 & 0\\
%     %& 0 & 0 & 1 & 1 & 1 & 0 & 1 & 0\\
% \end{tabular}\\
% Convertissez en décimal les deux nombres binaires suivant :
% 0001 0001 et 1011 1100 et 1001 1010. Vous pouvez vous aider du tableau donné ci-dessus.


% \begin{tabular}{lcccc}
%     & $16^3$ & $16^2$ & $16^1$ & $16^0$\\
%    & 4096 & $256$ & $16$ & $1$\\
%     \hline
%     $56_{(10)} = $ & 0 & 0 & 3 & 8\\
%     $12_{(10)} = $ & 0 & 0 & 0 & B\\
%  \end{tabular}\\

% Convertissez en décimal les nombres $2366_{(16)}$ , $A2EC_{(16)}$,  $AB_{(16)}$. Vous pouvez vous aider du tableau donnée ci-dessous.

% Trouver la représentation en base 16 puis en base 2 des nombres $112_{(10)}$ $8423_{(10)}$ \\

% D'après le tableau ASCII donné en annexe, décoder le mot suivant donné en hexa:
% carpe diem.

% Le codage d’une couleur se fait à l'aide de trois nombres compris chacun, en écriture décimale, entre 0 et 255 (code RVB). La couleur « Lie de vin » est codée, en écriture décimale, par \((172, 30, 68)\).

% Donner le codage hexadécimal correspondant à cette couleur.

% Que ce passe-t-il si je multiplie le nombre binaire 00110101 par 2? Si vous n'avez pas d'idée, convertissez ce nombre en décimal, multipliez par 2 puis remettez le en binaire.

% \end{document}
% \bfseries}
% \pagestyle{fancy}

% %\title{Thème 1 – Structures de données}
% %\date{}

% \begin{document}
% \lhead{Première NSI 2023-2024}
% \rhead{Thème 1 – Systèmes de numération}
% \rfoot{Chapitre 1}
% \fancyfoot[C]{\textbf{Page \thepage}}
% \input{page_de_garde}






% \chapter{Bases et changements de bases}


% Pour représenter l'information en machine, nous utilisons uniquement deux symboles dont l'une des raisons est qu'un appareil comme un ordinateur fonctionne avec des circuits électroniques. Il est donc plus facile de distinguer 2 états que plusieurs états qui peuvent être légèrement entachés d'erreur.

% Ces deux symboles ou valeurs (0 ou 1, vrai ou faux, ouvert ou fermé...) s'appellent des chiffres binaires ou bit en anglais (pour binary digit). Une variable qui ne peut prendre que 2 valeurs est également nommée booléen.

% C'est grâce à une multitude de petits circuits électroniques ne pouvant être chacun que dans 2 états, qu'une machine peut recevoir de l'information, la mémoriser, la modifier et la transmettre.

% \section{Les systèmes de numérations : représentation des entiers naturels}
% \subsection{Numération décimale :}
% La numération décimale utilise 10 chiffres : 0, 1, 2, 3, 4, 5, 6, 7, 8 et 9.\\
% L'écriture du nombre $275_{(10)}$ se traduit par :
% \[275= 2 \times 100 + 7 \times 10 + 5 \times 1= 2\times10^2 + 7\times10^1+ 5\times10^0\]
% On dit que $275$ est la représentation en base 10 appelée écriture décimale.\\

% Écrire une égalité semblable pour les nombres 1 234 et 208.\\
% \begin{tikzpicture}
% \draw[draw=black] (1,0) rectangle ++(18,2);
% \end{tikzpicture}


% \section{Base d'un système de numération}
% On appelle base d'un système de numération de position, le nombre de chiffres distincts permettant l'écriture de n'importe quel nombre entier naturel.\\

% Dans un système de base b, un nombre sera noté :
% \begin{equation*}
%     (a_n a_{n-1} ... a_1 a_0 )_b = a_0 \times b^0 + a_1 \times b^1 + ... + a_{n-1} \times b^*+++{n-1} + a_n \times b^n
% \end{equation*}

% où $a_i$ est un chiffre compris entre $0$ et $b-1$.

% \subsection{Décomposition en base 10}
% \begin{tabular}{lccccc}
% & $10^4$ & $10^3$ & $10^2$ & $10^1$ & $10^0$\\
% & $10000$ & $1000$ & $100$ & $10$ & $1$\\
% \hline
% $56_{(10)} = $ & 0 & 0 & 0 & 5 & 6\\
% $12_{(10)} = $ & 0 & 0 & 0 & 1 & 2\\

% \end{tabular}\\
% On déduit : \\
% $56_{(10)} = 5\times 10^1 + 6\times 10^0$\\
% $12_{(10)} = 1\times 10^1 + 2\times 10^0$\\


% Exemple pour le nombre $40121_{(10)}$ :\\
% \begin{tikzpicture}
%     \draw[draw=black] (1,0) rectangle ++(18,3);
% \end{tikzpicture}
% %(1234)10 = 4 x 100 + 3 x 101 + 2 x 102 + 1 x 103


% %
% En informatique, nous utilisons principalement 3 bases :
% \begin{itemize}
%     \item Le système binaire utilisant la base 2 : il est un concept essentiel de l'informatique. En effet, les processeurs des ordinateurs actuels sont composés de transistors ne gérant chacun que deux états. Un calcul informatique n'est donc qu'une suite d'opérations sur des paquets de 0 et de 1, appelés octets lorsqu'ils sont regroupés par huit (1 octet = 8 bits).

%     \item Le système décimal utilisant la base 10 : il est de nos jours le plus utilisé pour représenter des nombres, elle découle d'un choix naturel, dicté par le nombre de doigts des 2 mains.

%     \item Le système hexadécimal utilisant la base 16 : il est fréquemment utilisé en informatique car il est permet un compromis entre le code binaire des machines et une base de numération pratique à utiliser pour les ingénieurs (chaque chiffre correspond exactement à 4 chiffres binaires - 4 bits).

% \end{itemize}


% \subsection{Décomposition en base 2 (exemples) :}

% C'est donc le système de base 2 qu'utilisent les ordinateurs. Les chiffres sont au nombre de 2 : 0 et 1. Ainsi le nombre $12 = 1100_{(2)}$ est stocké par $4$ circuits mémoires respectivement dans les états $1, 1, 0, 0$.\\*

% \begin{tabular}{lcccccccc}
%     & $2^7$ & $2^6$ & $2^5$ & $2^4$ & $2^3$ & $2^2$ & $2^1$ & $2^0$\\
%     & $128$ & $64$ & $32$ & $16$ & $8$ & $4$ & $2$ & $1$\\
%     \hline
%     $12_{(10)} = $ & 0 & 0 & 0 & 0 & 1 & 1 & 0 & 0\\
%     $56_{(10)} = $ & 0 & 0 & 1 & 1 & 1 & 0 & 0 & 0\\
%     \end{tabular}\\
%     On déduit : \\
%     $56_{(10)} = 1\times 2^5 + 1\times 2^4 + 1\times 2^3 + 0\times 2^2 + 0\times 2^1 + 0\times 2^0 = 111000_{(2)}$\\
%     $12_{(10)} = 1\times 2^3 + 1\times 2^2 + 0\times 2^1 + 0\times 2^0 = 1100_{(2)}$\\

% Exprimer en base 2 les nombres 3, 6, et 14 pour illustrer cette réponse.\\
% \begin{tikzpicture}
%     \draw[draw=black] (1,0) rectangle ++(18,3);
% \end{tikzpicture}
    
% Un nombre binaire de huit chiffres est un \textbf{octet}.\\
% Quel est le plus grand nombre que l'on peut représenter avec un octet ?\\
% \begin{tikzpicture}
%     \draw[draw=black] (1,0) rectangle ++(18,3);
% \end{tikzpicture}

% \subsection{Numération hexadécimale :}
% La numération hexadécimale utilise 16 chiffres : 0, 1, 2, 3, 4, 5, 6, 7, 8, 9, A, B, C, D, E et F.\\
% A représente donc le nombre $10$, B le nombre $11$ etc ...\\
% \begin{tabular}{lccc}
%     & $16^2$ & $16^1$ & $16^0$\\
%     & $256$ & $16$ & $1$\\
%     \hline
%     $56_{(10)} = $ & 0 & 3 & 8\\
%     $12_{(10)} = $ & 0 & 0 & B\\
%     \end{tabular}\\
%     On déduit : \\
%     $56_{(10)} = 3\times 16^1 + 8\times 16^0 = 38_{(16)}$\\
%     $12_{(10)} = 12\times 16^0 = B_{(16)}$\\

% Traduire en nombre décimal $B8C_{(16)}$:\\
% \begin{tikzpicture}
%     \draw[draw=black] (1,0) rectangle ++(18,3);
% \end{tikzpicture}

% A quel nombre décimal correspond le nombre $FF_{(16)}$?\\
% \begin{tikzpicture}
%     \draw[draw=black] (1,0) rectangle ++(18,3);
% \end{tikzpicture}

% \section{Changement de base}
% \subsection{Passage de la base b à la base 10}
% Passer de l'écriture en base b d'un nombre à son écriture décimale est aisé, il suffit de faire la somme de la multiplication de chaque chiffre d'indice k par son poids correspondant $b^k$ (b étant la base d'origine).\\

% Exemple d'un nombre en hexadécimal vers un nombre binaire :
% Au préalable, on prendra soin de transcoder les lettres A à F en décimal (de 10 à 15) :
% \begin{itemize}
%     \item A vaut 10
%     \item B vaut 11
%     \item C vaut 12
%     \item D vaut 13
%     \item E vaut 14
%     \item F vaut 15.
% \end{itemize}
% Ainsi : $FA9_{(16)}= 15 \times 16^2 + 10 \times 16^1 + 9 \times 16^0 = 3840 + 160 + 9 = 4009_{(10)}$\\

% Donner l'écriture décimale de $76401_{(8)}$.\\
% \begin{tikzpicture}
%     \draw[draw=black] (1,0) rectangle ++(18,3);
% \end{tikzpicture}

% \subsection{Passage de la base 10 à la base b}
% Pour écrire les entiers naturels en base b, on a besoin de b chiffres. Quand on a n objets, on les groupe par paquets de b, qu'on regroupe à leur tour en paquets de b paquets, etc. Autrement dit, on fait une succession de divisions par b, jusqu'à obtenir un quotient égal à 1. Ainsi, si on veut trouver un résultat en binaire, on va diviser par 2. Si on veut de l'hexadécimal, on va diviser par 16.\\
% Un exemple pour le nombre $77_{(10)}$\\
% \begin{center}
%     
% \end{center}
% Le résultat sera donc $77_{(10)} = 1001101_{(2)}$ (lecture de bas vers le haut des restes)\\

% Trouver la représentation en base 16 puis en base 2 du nombre 6725 (écrits en base 10):\\
% \begin{tikzpicture}
%     \draw[draw=black] (1,0) rectangle ++(18,6);
% \end{tikzpicture}

% \subsection{Passerelle entre la base 2 et la base 16}
% Pour passer de la base 2 à la base 16 :
% \begin{itemize}
%     \item On décompose ce nombre par tranches de 4 bits à partir du bit de poids faible ( $2^4= 16$)
%     \item On complète la dernière tranche par des 0 s'il y a lieu.
%     \item On convertit chaque tranche en son symbole de la base 16.
%     \item On réécrit à sa place le nouveau symbole par changements successifs de chaque groupe de 4 bits.
% \end{itemize}
% Par exemple $111101_{(2)} = 3D_{(16)}$
% Pour passer de la base 16 à la base 2, cette conversion est l'opération inverse de la précédente.
% \begin{itemize}
%     \item On convertit chaque symbole hexadécimal du nombre en son écriture binaire (nécessitant au max 4 bits)
%     \item Pour chaque tranche de 4 bits, on complète les bits de poids fort par des 0 s'il y a
%     lieu.
%     \item On regroupe toutes les tranches de 4 bits à partir du bit de poids faible, sous
%     forme d'un seul nombre binaire.\\
% \end{itemize}

% Convertir $1100011010_{(2)}$ en base 16 et $7CF21_{(16)}$ en base 2.\\
% \begin{tikzpicture}
%     \draw[draw=black] (1,0) rectangle ++(18,6);
% \end{tikzpicture}


% \section{Contrôle}
% Comptez jusqu'à $1010_{(2)}$ en binaire.
% Comptez jusqu'à $F0_{(16)}$ en hexadécimal.
% De combien de bits est composé un octet ?
% Quel est le plus grand nombre qu'on peut représenter avec un octet ?
% Combien fait $FF_{(16)}+1_{(16)}$ ?
% \begin{tabular}{lcccccccc}
%     & $2^7$ & $2^6$ & $2^5$ & $2^4$ & $2^3$ & $2^2$ & $2^1$ & $2^0$\\
%     & $128$ & $64$ & $32$ & $16$ & $8$ & $4$ & $2$ & $1$\\
%     \hline
%      & 0 & 0 & 0 & 1 & 1 & 1 & 0 & 0\\
%     %& 0 & 0 & 1 & 1 & 1 & 0 & 1 & 0\\
% \end{tabular}\\
% Convertissez en décimal les deux nombres binaires suivant :
% 0001 0001 et 1011 1100 et 1001 1010. Vous pouvez vous aider du tableau donné ci-dessus.


% \begin{tabular}{lcccc}
%     & $16^3$ & $16^2$ & $16^1$ & $16^0$\\
%    & 4096 & $256$ & $16$ & $1$\\
%     \hline
%     $56_{(10)} = $ & 0 & 0 & 3 & 8\\
%     $12_{(10)} = $ & 0 & 0 & 0 & B\\
%  \end{tabular}\\

% Convertissez en décimal les nombres $2366_{(16)}$ , $A2EC_{(16)}$,  $AB_{(16)}$. Vous pouvez vous aider du tableau donnée ci-dessous.

% Trouver la représentation en base 16 puis en base 2 des nombres $112_{(10)}$ $8423_{(10)}$ \\

% D'après le tableau ASCII donné en annexe, décoder le mot suivant donné en hexa:
% carpe diem.

% Le codage d’une couleur se fait à l'aide de trois nombres compris chacun, en écriture décimale, entre 0 et 255 (code RVB). La couleur « Lie de vin » est codée, en écriture décimale, par \((172, 30, 68)\).

% Donner le codage hexadécimal correspondant à cette couleur.

% Que ce passe-t-il si je multiplie le nombre binaire 00110101 par 2? Si vous n'avez pas d'idée, convertissez ce nombre en décimal, multipliez par 2 puis remettez le en binaire.
